\chapterpage{controls}{设置视图并标记一个读数}
\begin{enumerate}
\item 切到 \ui{P-rho Diagram}，在 \ui{Plot Settings} 中将压强单位设为 Mbar。Hicks 表 I 对第 29400 号实验推得的石英密度为 6.03 g/cm³，压强为 4.79 Mbar。
\item 将密度坐标设为 5.4–6.8 g/cm³，压强坐标设为 3.5–6.3 Mbar，点击 \ui{Apply}，可以看清实验状态附近的两条曲线。
\item 在散点面板的手工输入区填入 6.030 和 4.790，点击 \ui{Add}。这样可以在密度–压强图上直接将论文推得的状态与 Hamel 石英曲线比较。
\end{enumerate}
\pic{settings}{散点标示 Hicks 论文给出的状态，两种标准材料的拟合曲线采用相同的坐标和单位。}{88mm}
\subhead{应当保存哪种标记}
记录目测读数时，把页签、两个量的名称和单位与手工点一起保留。记录计算交点时，可在散点区选用交点 ID。只要状态量齐全，程序也可以在其他相关图上标出该状态。要保留散点列表可保存工作流，若同事只需要阅读结果则可导出图像。
